% =============================================================================
% Chapter 18: Conclusion -- Part III (Extended Validation) and Overall Conclusion
% =============================================================================
\mychapter{18. Conclusion}

This part revisited three design choices made without controlled comparison earlier in the thesis,
loss function, preprocessing, and sample-boundary detection, and reported two further analyses,
network interpretability and a strengthened external validation, on the same PODFAM pipeline validated
in Part II. A controlled four-way loss-function comparison found BCE, not the combined Dice-Focal
formulation stated as primary in Part II, achieves the strongest validation Dice and IoU, and further
established that Part II's own headline generalisation result was already trained with BCE, correcting
rather than merely refining Part II's stated choice. A held-out test sample, evaluated across every
configuration tested in this part, consistently scored far below its corresponding validation sample,
0.12--0.21 Dice against 0.70--0.73, replicating, on an independent split and set of configurations, the
same pattern Part II first reported for its own single trained model (Section~10.3). This
generalisation gap, not any single preprocessing or masking result, is this part's central finding: it
is a property of how few distinct samples this pipeline has to work with, not of any one training run,
and it remains the pipeline's primary open problem going into any future work on this dataset.

Within that constraint, beam-hardening/ring-artefact correction and a shape-agnostic replacement for
the circular sample-boundary assumption each moved held-out test Dice in the right direction (0.1218
$\to$ 0.1693 $\to$ 0.2119), evidence that these specific engineering changes improve genuine
generalisation rather than only fitting the validation sample more closely. The shape-aware boundary
method also corrected a two-order-of-magnitude artefact in sample\_06's reported defect fraction
(2.92\% $\to$ 0.030\%) and recovered the same external-validation accuracy as the circular mask it
replaced (Table~\ref{tab:mask_r_comparison}) while removing an assumption already known not to hold for
two of the seven PODFAM samples. A parallel investigation of Canny edge detection as a further
alternative is reported as a negative result: it independently corroborates the circular-fit method on
samples where both apply, but exhibits a confirmed, data-driven failure mode on the non-circular
samples that the adopted shape-aware method does not, and is documented here specifically so that
choice is not revisited without this evidence.

A gradient-based attribution analysis of the trained detector, run on the verified BCE checkpoint,
found that the three encoder blocks together account for roughly two-thirds of the network's
attributed defect-discriminating signal, with the block closest to the output the least influential of
all nine, offering a concrete, evidence-based direction (strengthening early-encoder feature extraction
or its skip connections) for future architecture work on this task. Separately, checking Kim et al.'s
own reference porosity values against their paper's stated methodology found they were themselves
cross-validated against an independent gravimetric method, strengthening the interpretation of this
thesis's $r = 0.995$ external-validation result as agreement with an already physically-grounded
reference rather than an XCT-internal comparison alone. A further, still-unconfirmed observation, that
the two NIST samples this thesis excludes from training were scanned at a substantially finer voxel
size than the rest of the set, is reported as a plausible explanation worth testing directly in any
continuation of this work, rather than as a settled finding.

Taken together with Parts I and II, this document now covers method development on a public benchmark
(NIST CoCr, Part I), full-volume classical-versus-deep-learning validation on real industrial
process-monitoring data (PODFAM, Part II), and a controlled revisiting of that validation's own design
choices alongside two further, independent lines of evidence, interpretability and external-reference
verification (Part III). The overall picture that emerges is consistent across all three parts: the
pipeline's classical thresholding stage, particularly Bernsen local adaptive thresholding, is validated
against independently published, physically-grounded reference data with high confidence; the trained
detector built on top of it produces plausible, well-motivated predictions within the sample or samples
it was trained and validated on, but its performance on data it has never seen falls far short of this
thesis's stated acceptance criteria in every configuration tested; and every major result in this
document, favourable or not, was reached only after an explicit, reported cross-check, itself a
recurring finding of this work, that a classical method's or a trained model's output looking plausible
in isolation is not sufficient grounds to trust it without independent verification.
